% !TEX root = ../print.tex
% Draft \S8 --- Examples and moments
%
% Source: draft/hemigroup-causal-scale-space-kernels.md, lines 254-278.
% Chapter 8 of the blueprint; the shared counter puts every number here on the draft's own.
%
% The draft calls 8.1-8.3 Examples. They are transcribed as PROPOSITIONS, because they are
% load-bearing rather than illustrative: \ref{thm:locality} classifies the Bessel-K family as
% the order-2 local case, \ref{prop:local-ladder} builds on the Gamma family, and \S\S13-14
% (outside this transcription) rest on both. A \uses{} edge must point at a node, and an
% `example` is not one -- linkage does not count it a statement environment and the dependency
% graph does not carry it.
%
% One node per family, not a definition/proposition pair. A pair would occupy 8.1-8.6 and push
% Proposition 8.4 to 8.7, breaking the numbering agreement with the draft for the sake of a
% distinction the content does not need: each of these states "this F is admissible and its
% kernels are these", which is one claim. Same reasoning as \ref{prop:bernstein-toolbox}.
%
% Two [A] nodes here, the first since chapter 2: \ref{prop:bessel-family} on ledger A8
% (Halgreen) and \ref{prop:moment-criterion} on A7 (Sato). Both say more than their citation carries,
% so both Assignment clauses divide the statement explicitly.

Throughout this chapter we work in the canonical gauge, so that
$\phi_x = \Lap^{-1}[e^{-F(x s)}]$, and $T_x$ denotes a random variable with density $\phi_x$.

% draft: Example 8.1
\begin{proposition}[the extremal stable family; the 2005 kernels]\label{prop:stable-family}
\uses{lem:selfdecomposable-exponents,cor:semigroup-case,prop:moment-criterion}
For $0 < \alpha < 1$ the exponent $F(s) = s^\alpha$ is of the form (7.1), with Lévy density
factor
\[
  k(t) = \frac{\alpha}{\Gamma(1-\alpha)}\, t^{-\alpha},
\]
which is nonincreasing. Every moment of $T_x$ is infinite, so the delay cannot be measured by
the mean; measured by the mode it scales as $x^{1/\alpha}$.
\statusT
\end{proposition}

\begin{proof}
$s^\alpha = \frac{\alpha}{\Gamma(1-\alpha)}\int_0^\infty (1 - e^{-st})\,t^{-\alpha-1}\,dt$ is
the standard integral representation, which is (7.1) with $b_0 = 0$ and
$k(t)/t = \alpha t^{-\alpha-1}/\Gamma(1-\alpha)$; $k$ is a negative power, hence
nonincreasing. Since $\int_1^\infty t^{n-1}k(t)\,dt = \infty$ already for $n = 1$, all moments
are infinite by \ref{prop:moment-criterion}. The mode scaling is the self-similarity
$T_x \stackrel{d}{=} x^{1/\alpha} T_1$ of the $\alpha$-stable subordinator, which is the case
$\chi(x) = x^{1/\alpha}$ of \ref{cor:semigroup-case}.
\end{proof}

% draft: Example 8.2
\begin{proposition}[the Gamma family]\label{prop:gamma-family}
\uses{lem:selfdecomposable-exponents,thm:main-characterization,rem:sato-reading}
For $\gamma > 0$ the exponent $F(s) = \gamma\,\log(1 + s)$ is of the form (7.1) with
$k(t) = \gamma\, e^{-t}$, which is nonincreasing. Its kernels are
\[
  e^{-F(xs)} = (1 + x s)^{-\gamma}, \qquad
  \phi_x(t) = \frac{t^{\gamma - 1} e^{-t/x}}{\Gamma(\gamma)\, x^{\gamma}},
\]
the Gamma density of fixed shape $\gamma$ and scale $x$. All moments are finite, with
$\Ex\,T_x = \gamma x$ and $\operatorname{Var} T_x = \gamma x^2$. The increment measures
$\mu_{a,b}$ have Laplace transform $\bigl(\tfrac{1 + as}{1 + bs}\bigr)^{\gamma}$ and are the
increments of the Gamma-type Sato subordinator.
\statusT
\end{proposition}

\begin{proof}
The Frullani integral gives
$\log(1+s) = \int_0^\infty (1 - e^{-st})\,e^{-t}\,t^{-1}\,dt$, so $F$ has the form (7.1) with
$b_0 = 0$ and $k(t) = \gamma e^{-t}$, nonincreasing; hence $F$ is admissible by
\ref{lem:selfdecomposable-exponents}. Then $e^{-F(xs)} = (1+xs)^{-\gamma}$, which is the
Laplace transform of the stated Gamma density. The moments are those of a Gamma law, and all
are finite because $\int_1^\infty t^{n-1}e^{-t}dt < \infty$ for every $n$. The increment
transform is the ratio $e^{-F(bs)}/e^{-F(as)}$, and \ref{rem:sato-reading} identifies the
family as the marginals of a Sato subordinator.
\end{proof}

% draft: Example 8.3
\begin{proposition}[the Bessel-K family; preview of the memory line]\label{prop:bessel-family}
\uses{lem:selfdecomposable-exponents,thm:main-characterization,prop:moment-criterion}
For $a > 0$ let $T_1 \stackrel{d}{=} 1/(2\gamma_a)$ be the scaled inverse-gamma law, where
$\gamma_a$ has the Gamma density of shape $a$, and let $F$ be its Laplace exponent. Then
\[
  e^{-F(s)} = \frac{2^{1-a}}{\Gamma(a)}\,\bigl(\sqrt{2s}\bigr)^{a} K_a\bigl(\sqrt{2s}\bigr),
\]
and $T_1$ is self-decomposable, so $F$ is of the form (7.1) and the family is admissible. In
the parabolic gauge $\tilde x = x^2/2$ these are the boundary-hitting kernels of the
Bessel-type media $\tfrac12\partial_x^2 + \tfrac{\beta}{x}\partial_x$ with $a = \tfrac12 -
\beta$, and the moments satisfy $\Ex\,T_x^n < \infty$ if and only if $a > n$. This family is
the subject of the memory-line development, Theorems 10.4, 11.6 and 12.5.
\statusA\quad\ledger{A8}\quad\emph{(self-decomposability of the generalized inverse Gaussian
laws; \sscite{halgreen1979self}. A cited interface, not reproven here.
\textbf{Assignment.} \ledger{A8} carries exactly one clause: that $T_1$ is self-decomposable,
which is what makes $F$ admissible and puts the family inside
\ref{thm:main-characterization}. Everything else is \textbf{[T]} and proved below --- the
evaluation of the transform in terms of $K_a$ is a classical integral representation, and the
moment criterion is \ref{prop:moment-criterion} applied to this $k$. The parabolic-gauge reading is a
forward reference to Chapter~11 and is not asserted here. Note also that Halgreen proves the
stronger statement that the law is a generalized Gamma convolution; self-decomposability
follows through $\GGC \subset \SDclass$, which is a clause of \ledger{A16} --- so A8 and A16
rest on the same tower and should be reviewed together.)}
\end{proposition}

\begin{proof}
Write $\gamma_a$ for a Gamma$(a,1)$ variable and $T_1 = 1/(2\gamma_a)$. Substituting
$u = 1/(2t)$ in the Gamma density gives the density of $T_1$,
\[
  \phi_1(t) = \frac{t^{-a-1}\,e^{-1/(2t)}}{2^{a}\,\Gamma(a)}, \qquad t > 0 .
\]
The classical integral representation
$\int_0^\infty t^{\nu-1}\exp(-A/t - Bt)\,dt = 2\,(A/B)^{\nu/2}K_\nu\bigl(2\sqrt{AB}\bigr)$,
taken at $\nu = -a$, $A = \tfrac12$, $B = s$, together with $K_{-a} = K_a$, gives
\[
  \Ex\,e^{-sT_1}
  = \frac{1}{2^{a}\Gamma(a)}\int_0^\infty t^{-a-1}e^{-st - 1/(2t)}\,dt
  = \frac{2\,(2s)^{a/2}\,K_a\bigl(\sqrt{2s}\bigr)}{2^{a}\Gamma(a)}
  = \frac{2^{1-a}}{\Gamma(a)}\bigl(\sqrt{2s}\bigr)^{a}K_a\bigl(\sqrt{2s}\bigr).
\]
Self-decomposability of $T_1$ is \ledger{A8}; by
\ref{lem:selfdecomposable-exponents} its exponent is therefore of the form (7.1), and
\ref{thm:main-characterization} makes the family admissible. Finally, the Lévy density factor
of this $F$ satisfies $k(t) \sim a$ as $t \to 0$ and decays like $t^{-a}$ at infinity, so
$\int_1^\infty t^{n-1}k(t)\,dt < \infty$ exactly when $n < a$; \ref{prop:moment-criterion} converts
that into $\Ex\,T_x^n < \infty$ iff $a > n$.
\end{proof}

% draft: Proposition 8.4
\begin{proposition}[the mean delay and the influence curve]\label{prop:moments}
\uses{lem:selfdecomposable-exponents,thm:main-characterization}
$\Ex\,T_x = x\,F'(\zp) = x\bigl(b_0 + \int_0^\infty k(t)\,dt\bigr)$, finite if and only if $k$
is integrable at infinity. When the mean is finite, the influence curve of
\sscite{fagerstrom2005temporal}, \S7, may be defined by the mean delay, and it is exactly
linear in the canonical gauge:
\[
  t_0 - t_m = F'(\zp)\; x .
\]
\statusT\quad\emph{(monotone convergence in (7.1) and nothing else --- no citation is needed and
none is made. The higher-moment criterion, which \emph{does} need one, was split off as
\ref{prop:moment-criterion}; before the split this node reported that entry's cost for a result
that does not incur it. What is \emph{not} claimed here is anything about the mode, which is the
delay measure the stable family forces one to use --- see \ref{prop:stable-family}.)}
\end{proposition}

\begin{proof}
Differentiating the transform at the origin,
$\Ex\,T_x = -\partial_s e^{-F(xs)}\big|_{s=\zp} = x F'(\zp)$, and monotone convergence in
(7.1) gives $F'(\zp) = b_0 + \int_0^\infty k(t)\,dt$, finite exactly when $k$ is integrable at
infinity. The influence curve is then linear in $x$ because $\Ex\,T_x$ is.
\end{proof}

% draft: Remark 8.5
\begin{remark}[heavy tails are the semigroup's doing]\label{rem:heavy-tails}
\uses{thm:main-characterization,cor:semigroup-case,prop:moment-criterion,prop:gamma-family}
The heavy tails of the 2005 kernels are an artifact of the semigroup axiom, not of causality
together with covariance. Within \ref{thm:main-characterization} the semigroup sub-case
(\ref{cor:semigroup-case}) is precisely the subset with $k$ a pure power --- the only $k$
compatible with homogeneity --- and every non-power nonincreasing $k$, for instance
$k(t) = \gamma e^{-t}$, yields a causal, continuously scale-covariant family with finite
moments of all orders.
\end{remark}

% additive: not in the draft. The higher-moment half of Proposition 8.4, split off so that the
% mean-delay identity -- which needs no citation -- does not report ledger A7's cost. Numbered
% 8.6, a number the draft does not use, its Section 8 ending at 8.5.
\begin{proposition}[the moment criterion]\label{prop:moment-criterion}
\uses{prop:moments,lem:selfdecomposable-exponents,thm:main-characterization}
$\Ex\,T_x^{\,n} < \infty$ if and only if $\int_1^\infty t^{\,n-1} k(t)\,dt < \infty$.
\statusA\quad\ledger{A7}\quad\emph{(the moment criterion for infinitely divisible laws;
\sscite{sato1999levy}, Thm.~25.3. A cited interface, not reproven here.
\textbf{Assignment.} \ledger{A7} carries the equivalence itself and nothing else --- that a
moment of the law is finite exactly when the corresponding integral against the Lévy measure
over $\{t > 1\}$ is finite. Translating that integral into the stated one is \textbf{[T]}: the
Lévy measure of $T_x$ is $\nu(dt) = k(t)t^{-1}dt$ up to the scaling by $x$, so
$\int_{t>1}t^n\,\nu(dt) = \int_1^\infty t^{\,n-1}k(t)\,dt$. The $n = 1$ case is \emph{not}
carried by the ledger --- it is \ref{prop:moments}, proved by monotone convergence --- which is
the reason for the split.)}
\end{proposition}

\begin{proof}
The Lévy measure of $T_x$ is $\nu(dt) = k(t)t^{-1}dt$ up to the scaling by $x$, so
$\int_{t>1} t^n\,\nu(dt) = \int_1^\infty t^{\,n-1}k(t)\,dt$; \ledger{A7} equates finiteness of
that integral with finiteness of $\Ex\,T_x^n$.
\end{proof}
